\begin{table}[t]
\centering
\caption{实验一的端到端系统行为（片段或遮挡过程的中位数）。所有指标越低越好，粗体为每行最优；原始片段值与 IQR 见图~\ref{fig:exp1-final}。}
\label{tab:exp1-performance}
\normalsize
\setlength{\tabcolsep}{1.5pt}
\renewcommand{\arraystretch}{1.08}
\begin{tabularx}{\columnwidth}{@{}>{\raggedright\arraybackslash}Xcccc@{}}
\toprule
指标 $\downarrow$ & Arrival & Capture & One-Euro & EgoAnchor \\
\midrule
\multicolumn{5}{@{}l}{\textit{静止与遮挡}} \\
头动泄漏 P95 (mm) & 43.67 & 13.03 & 10.65 & \textbf{0.82} \\
绝对注册 P95 (mm) & 44.31 & 17.54 & 14.00 & \textbf{6.60} \\
\shortstack[l]{静止帧间增量\\P95 (mm)} & 11.70 & 4.30 & 0.85 & \textbf{0.06} \\
遮挡平移 P95 (mm) & 18.23 & 16.93 & 10.41 & \textbf{4.85} \\
Start-transition (ms) & \textbf{157.04} & 263.76 & 333.98 & 510.36 \\
\midrule
\multicolumn{5}{@{}l}{\textit{持续平移}} \\
有效时延 (ms) & \textbf{202.50} & 255.00 & 380.00 & 360.00 \\
对齐 RMSE (mm) & 18.57 & 17.14 & 15.69 & \textbf{9.12} \\
\shortstack[l]{当前时刻\\RMSE (mm)} & \textbf{84.74} & 103.63 & 117.61 & 125.83 \\
\midrule
\multicolumn{5}{@{}l}{\textit{持续旋转}} \\
有效时延 (ms) & \textbf{255.00} & 257.50 & 385.00 & 345.00 \\
对齐 RMSE (deg) & 6.09 & 5.84 & 5.55 & \textbf{4.81} \\
\shortstack[l]{当前时刻\\RMSE (deg)} & 25.44 & \textbf{25.35} & 34.10 & 31.88 \\
\bottomrule
\end{tabularx}
\end{table}
